% Set up the document
\documentclass[a4paper, 12pt, oneside]{GuThesis}
\usepackage{GuThesis}

\begin{document}

% Tez Türü ----------------------------------------------------
\thesistype{0}	% Tasarım = 0, Bitirme = 1
% -------------------------------------------------------------
% Tez Başlığı -------------------------------------------------
\thesistitle{TEZ BAŞLIĞI}
\thesistitlesmall{Tez Başlığı}
\thesistitleeng{THESIS NAME}
\thesistitlesmalleng{Thesis Name}
% -------------------------------------------------------------
% Tez Tarihi --------------------------------------------------
\thesisdate{OCAK}{2022}
\thesisdateeng{JANUARY}{2022}
\thesisdefdate{12/06/2022}
% -------------------------------------------------------------
% Yazar Bilgileri ---------------------------------------------
% Tez grubu kaç kişi ise o kadar yazarı aktif edin.
% Yazar 1
\thesisautoroneen{1} % Yazar Pasif = 0, Yazar Aktif = 1
\thesisauthorone{ADI SOYADI 1}{Numarası}
\thesisauthoronesmall{Adı SOYADI}
% Yazar 2
\thesisautortwoen{1} % Yazar Pasif = 0, Yazar Aktif = 1
\thesisauthortwo{ADI SOYADI 2}{Numarası}
\thesisauthortwosmall{Adı SOYADI}
% Yazar 3
\thesisautortreeen{0} % Yazar Pasif = 0, Yazar Aktif = 1
\thesisauthortree{ADI SOYADI 3}{Numarası}
\thesisauthortreesmall{Adı SOYADI}
% Yazar 4
\thesisautorfouren{0} % Yazar Pasif = 0, Yazar Aktif = 1
\thesisauthorfour{ADI SOYADI 4}{Numarası}
\thesisauthorfoursmall{Adı SOYADI}
% -------------------------------------------------------------
% Danışman ve Jüri Üyeleri ------------------------------------
% Tez danışmanı
\supapbig{Prof. Dr. ADI SOYADI}
\supap{Prof. Dr. Adı SOYADI}
% Jüri üyeleri
% Jüri 1
\juryaponebig{Prof. Dr. ADI SOYADI 1}
\juryapone{Prof. Dr. Adı SOYADI 1}
% Jüri 2
\juryaptwobig{Prof. Dr. ADI SOYADI 2}
\juryaptwo{Prof. Dr. Adı SOYADI 2}
% Jüri 3
\jurythreeen{0} % Jüri Pasif = 0, Jüri Aktif = 1
\juryapthreebig{Prof. Dr. ADI SOYADI 3}
\juryapthree{Prof. Dr. Adı SOYADI 3}
% Bölüm başkanı
\dephead{Doç. Dr. Mustafa Engin BAŞOĞLU}
% -------------------------------------------------------------
% Tez Sayfa Sayısı --------------------------------------------
\dummypages{20}
% -------------------------------------------------------------

% Özet Sayfası Bilgileri --------------------------------------
% Türkçe anahtar kelimeler
\turabkeys{Zamanlama analizi, İstatiksel modelleme, Benzetim.}
% İngilizce anahtar kelimeler
\engabkeys{Time analysis, Statistical modelling, Simulation.}
% -------------------------------------------------------------

% Sayfa Numaraları Düzenlemesi (Değiştirmeyiniz) --------------
\maketitlex

\clearpage
\BgThispage
\thispagestyle{empty}
\secpage
\thispagestyle{empty}
\setstretch{1}
\approform
\setstretch{1.5}
\pagenumbering{Roman}
\setcounter{page}{3}
% -------------------------------------------------------------

% Önsöz -------------------------------------------------------
\addcontentsline{toc}{section}{ÖNSÖZ}
\begin{preface} \par 
İlk sayfa niteliğinde yazılır ve bir sayfayı geçmez. Tezi hazırlayanın belirtmek istediği özel mesajları, kişisel görüş, amaç ve dileklerini kapsar. Tezi destekleyen kuruluşlar varsa, bunlardan söz edilebilir ve istenirse, ilgililere teşekkür edilir. Ayrıca teşekkür sayfası bulunmaz. Sayfanın üst kısmına, ortaya gelecek şekilde büyük harflerle (kalın ve koyu karakter) ÖNSÖZ yazılır. Önsözün sağ alt kısmında yazarın ismi yazılır. Yazar adının hemen altına tezin teslim edildiği il ve yıl yazılır.
\end{preface}
% -------------------------------------------------------------
% Tez etik beyannamesi ----------------------------------------
\addcontentsline{toc}{section}{TEZ ETİK BEYANNAMESİ}
\thesdec
% -------------------------------------------------------------
% İçindekiler -------------------------------------------------
\addcontentsline{toc}{section}{İÇİNDEKİLER}
\tableofcontents
\newpage
% -------------------------------------------------------------
% Türkçe özet -------------------------------------------------
\addcontentsline{toc}{section}{ÖZET}
\begin{turabstract}

Özet hazırlanırken 1.5 satır boşluk bırakılır. Tezlerde özet/abstract 300 kelimeden az olmamak kaydıyla 2 sayfayı geçmemelidir,  Özetlerde tezde ele alınan konu kısaca tanıtılarak, kullanılan yöntemler ve ulaşılan sonuçlar belirtilir.  Özetlerde kaynak, şekil, çizelge ve dipnot kullanılmamalıdır.  ÖZET  birinci dereceden başlık formatında (önce 72, sonra 18 punto aralık bırakılarak ve 1.5 satır aralıklı olarak) yazılmalıdır. ÖZET'in altına tezin türü (Yüksek Lisans veya Doktora)  belirtildikten sonra büyük harflerle sayfa ortalanarak (büyük harflerle) tezin başlığı yazılmaldır.

\end{turabstract}
% -------------------------------------------------------------
% İngilizce özet ----------------------------------------------
\addcontentsline{toc}{section}{SUMMARY}
\begin{engabstract}

1.5 line spacing must be set for Abstract. The abstract must have 300 words minimum and span 2 pages. A summary must briefly mention the subject of the thesis, the method(s) used and the conclusions derived. References, figures and tables must not be given in Summary. Below the Abstract, the thesis title in first level title format with capital letters (i.e., 72 pt before and 18 pt after paragraph spacing, and 1.5 line spacing) must be placed. Below the title, the expression must be written horizontally centered.

\end{engabstract}
% -------------------------------------------------------------
% Şekiller dizini ---------------------------------------------
\addcontentsline{toc}{section}{ŞEKİLLER DİZİNİ}
\addtocontents{lof}{~\hfill\underline{\textbf{Sayfa No}}\par}
\setstretch{1}
\listoffigures
\setstretch{1.5}
\vfill
\newpage
% -------------------------------------------------------------
% Tablolar dizini ---------------------------------------------
\addcontentsline{toc}{section}{TABLOLAR DİZİNİ}
\addtocontents{lot}{~\hfill\underline{\textbf{Sayfa No}}\par}
\setstretch{1}
\listoftables
\setstretch{1.5}
\vfill
\newpage
% -------------------------------------------------------------
% Semboller dizini --------------------------------------------
\addcontentsline{toc}{section}{SEMBOLLER DİZİNİ}
\section*{\centerline{SEMBOLLER DİZİNİ}}
\begin{center}
    
\begin{tabular}{ l l }
 \textbf{CFL} & : Courant-Friedrichs-Lewy kriteri \\ 
 \textbf{EOS} & : Hal denklemi (Equation of state) \\ 
 \textbf{FCT} & : Akı-Düzeltmeli Taşınım Algoritması (Flux- Corrected Transport) \\ 
 \textbf{HB} & : Hidrojen bağları (hydrogen bonds) \\ 
 \textbf{LCPFCT} & : Laboratory for Computational Physics, Flux- Corrected Transport \\ 
 \textbf{PIV} & : Parçacık Görüntü Hızölçer (Particle Image Velocimetry)  \\
 \parbox[][][c]{80pt}{\,} & 
\end{tabular}

\end{center}
\vfill
\newpage
% -------------------------------------------------------------

\pagenumbering{arabic}{

\chapter{GENEL BİLGİLER}

\section{Giriş}

\par Metin metin metin metin metin metin metin metin metin metin metin metin metin metin metin metin metin metin metin metin metin metin metin metin metin metin metin metin metin metin metin metin metin metin metin metin metin metin metin metin metin metin metin metin metin metin metin.

\section{Alt Başlık (Düzey 1)}

Metin metin metin metin metin metin metin metin metin metin metin metin metin metin metin metin metin metin metin metin metin metin metin metin metin metin metin metin metin metin metin metin metin metin metin metin metin metin metin metin metin metin metin metin metin metin metin metin metin metin metin metin metin metin metin metin metin metin metin metin metin metin metin metin metin metin metin metin metin metin metin metin metin metin metin metin metin metin metin metin metin metin metin metin metin.

\begin{figure}[h]
    \centering
    \includegraphics{csk_blok_diyagram.png}
    \caption{Renk kaydırmalı anahtarlama modülasyonu blok diyagramı}
\end{figure}

Metin metin metin metin metin metin metin metin metin metin metin metin metin metin metin metin metin metin metin metin metin metin metin metin metin metin metin metin metin metin metin metin metin metin metin metin metin metin metin metin metin metin metin metin metin metin metin metin metin metin metin metin metin metin metin

\subsection{Alt Başlık (Düzey 2)}

Metin metin metin metin metin metin metin metin metin metin metin metin metin metin metin metin metin metin metin metin metin metin metin metin metin metin metin metin metin metin metin metin metin metin metin metin metin metin metin metin metin metin metin metin metin metin metin metin metin metin metin metin metin metin metin metin metin.

\section{Alt Başlık (Düzey 1)}

Metin

\subsection{Alt Başlık (Düzey 2)}

Metin

\subsubsection{Alt Başlık (Düzey 3)}

Metin

\paragraph{Alt Başlık (Düzey 4)}

Metin metin metin metin metin metin metin metin metin metin metin metin metin metin metin metin metin metin metin metin metin metin metin metin metin metin metin metin metin metin metin metin metin metin metin metin metin metin metin metin metin metin metin metin metin metin.

\chapter{YAPILAN ÇALIŞMALAR}

\section{Giriş}

Metin metin metin metin metin metin metin metin metin metin metin metin metin metin metin metin metin metin metin metin metin metin metin metin metin metin metin metin metin metin metin metin metin metin metin metin metin metin metin metin metin metin metin metin metin metin metin metin metin metin metin metin metin metin metin metin metin metin metin metin metin.

\begin{figure}[h]
    \centering
    \includegraphics[width=0.8\linewidth]{color_space.png}
    \caption{Standartta belirtilen (110, 001, 000) renk bandı kombinasyonu}
\end{figure}

Metin metin metin metin metin metin metin metin metin metin metin metin metin metin metin metin metin metin metin metin metin metin metin metin metin metin metin metin metin metin metin metin metin metin metin metin metin metin metin metin metin metin metin metin metin metin metin metin metin metin metin metin metin metin metin metin metin metin metin metin metin metin metin metin metin metin metin metin metin metin metin metin metin metin metin metin metin metin metin metin metin metin metin metin metin metin metin metin metin metin metin metin metin metin metin metin metin metin metin metin metin metin metin metin

\section{Alt Başlık (Düzey 1)}

Metin metin metin metin metin metin metin metin metin metin metin metin metin metin metin metin metin metin metin metin metin metin metin metin metin metin metin metin metin metin metin metin metin metin metin metin metin metin metin metin metin metin metin metin metin metin.

\begin{table}[h]
\caption{Çözümün sayısal ağ yapısından bağımsızlığının araştırılması}
\centering
\begin{tabular}{m{0.1\linewidth} m{0.25\linewidth} m{0.5\linewidth}}
 \hline
 \textbf{Kanal} & \textbf{Walsh Kodu} & \textbf{Gönderilecek Eğitim Dizisi} \\ 
 \hline
 Kanal i & W(1,4) = 1, -1, 1, -1 & 1, 1, 0, 0, 1, 1, 0, 0 \\ 
 \hline
 Kanal j & W(2,4) = 1, 1, -1, -1 & 1, 1, 1, 1, 0, 0, 0, 0 \\
 \hline
 Kanal k & W(3,4) = 1, -1, -1, 1 & 1, 1, 0, 0, 0, 0, 1, 1 \\
 \hline
\end{tabular}
\end{table}

Metin metin metin metin metin metin metin metin metin metin metin metin metin metin metin metin metin metin metin metin metin metin metin metin metin metin metin metin metin metin metin metin metin metin metin metin metin metin metin metin metin metin metin metin metin metin metin metin metin metin metin metin metin metin metin metin metin metin metin.

\subsection{Alt Başlık (Düzey 2)}

Metin metin metin metin metin metin metin metin metin metin metin metin metin metin metin metin metin metin metin metin metin metin metin metin metin metin metin metin metin metin metin metin metin metin metin metin

Buna göre çoklu LED ve çoklu yansıma durumları için kanal tepkesi (\ref{ch_response}) numaralı bağıntıda verildiği gibi hesaplanır \cite{Reference1}.

\begin{equation}\label{ch_response}
    h(t) = \sum_{n=1}^{N_{LED}}\sum_{k=0}^\infty h^{(k)}(t, S_n)
\end{equation}

Burada $N_{LED}$ LED sayısını, $h^{(k)}$ ise $k$ adet yansıma durumunda oluşan kanal tepkesini temsil etmekte ve (2.2)’de verildiği gibi hesaplanmaktadır \cite{Reference2}.

\chapter{BULGULAR VE TARTIŞMALAR}

\section{Giriş}

Metin metin metin metin metin metin metin metin metin metin metin metin metin metin metin metin metin metin metin metin metin metin metin metin metin metin metin metin metin metin metin metin metin metin metin metin metin metin metin metin metin metin metin metin metin metin metin metin metin metin metin metin metin metin metin metin metin metin metin metin metin metin metin metin metin metin metin metin metin metin metin metin metin metin metin metin metin metin metin metin metin metin metin metin metin.

\begin{equation}\label{ch_response2}
    h(t) = \sum_{n=1}^{N_{LED}}\sum_{k=0}^\infty h^{(k)}(t, S_n)
\end{equation}

Metin metin metin metin metin metin metin metin metin.

\subsection{Alt Başlık (Düzey 2)}

Metin

\chapter{SONUÇLAR}

Metin metin metin metin metin metin metin metin metin metin metin metin metin metin metin metin metin metin metin metin metin metin metin metin metin metin metin metin metin metin metin metin metin metin metin metin metin metin metin metin metin metin metin metin metin metin metin metin metin metin metin metin metin metin metin metin metin metin metin metin metin metin metin metin metin metin metin metin metin metin metin metin metin metin metin metin metin metin metin metin metin metin metin metin metin metin metin metin metin metin metin metin metin metin metin metin metin metin metin metin metin.

\begin{table}[h]
    \caption{Örnek tablo}
    \centering
    \begin{tabular}{|m{0.3\linewidth}|m{0.3\linewidth}|m{0.3\linewidth}|}
         \hline
         & & \\
         \hline
         & & \\
         \hline
    \end{tabular}
\end{table}

Metin metin

\vfill
\newpage



\bibliographystyle{GuRefStyle}
\setstretch{1}
\bibliography{Referances.bib}
\setstretch{1.5}
\vfill
\newpage


\chapter{EKLER}
\par\bigskip

\textbf{Ek 1}
\begin{table}[htb]
\centering
\caption{Ekler bölümünde çizelge örneği}
\label{my-label}
\begin{tabular}{|l|l|l|l|}
\hline
Kolon1 & Kolon2 & Kolon3 & Kolon4 \\ \hline
Satır1 & Satır1 & Satır1 & Satır1 \\ \hline
Satır2 & Satır2 & Satır2 & Satır2 \\ \hline
Satır3 & Satır3 & Satır3 & Satır3 \\ \hline
\end{tabular}
\end{table}

\textbf{EK 2}
\begin{figure}[h]
    \centering
    \includegraphics{ek1.png}
    \caption{Eş merkezli silindirik borularda yerdeğiştiren akışkanın ısı transferi.}
\end{figure}

\vfill
\newpage

}

\addcontentsline{toc}{section}{ÖZGEÇMİŞ}
\chapter*{\centerline{ÖZGEÇMİŞ}}
Ad SOYAD, ............. tarihinde .......'da doğdu. İlk ve orta öğrenimini .............. İlköğretim Okulu'nda, lise öğrenimini ise ................ Lisesi'nde tamamladı. ..... yılında Gümüşhane Üniversitesi Mühendislik ve Doğa Bilimleri Fakültesi Elektrik Elektronik Mühendisliği Bölümü'nde lisans eğitimine başladı.


\end{document}
